\section*{Заключение}
\addcontentsline{toc}{section}{Заключение}

Проведенный комплексный историко-политологический и социологический анализ статуса, динамики и специфики взаимодействия научно-технической интеллигенции с институтами государственной власти позволяет сформулировать следующие итоговые выводы:

\begin{enumerate}
    \item Научно-техническая интеллигенция в постиндустриальном обществе трансформируется в специфическую макроструктуру --- «техноструктуру» (по Дж. Гэлбрейту), обладающую монополией на специальное знание, необходимое для государственного планирования \cite{Gelbreyt}. В системе политического анализа профессора А. И. Соловьева НТИ функционирует как влиятельная группа интересов, чья артикуляция требований реализуется через неокорпоративистские модели экспертного участия и лоббизма \cite{Solovyev}.
    \item Историческая ретроспектива доказывает этатистский характер генезиса отечественной НТИ от петровских модернизационных сдвигов до советского ВПК. Российское научно-техническое сословие исторически формировалось под патронажем государства, получая привилегии и высокий статус в обмен на лояльность и решение оборонных задач. Однако оборотной стороной этой модели всегда выступали жесткий идеологический контроль и политическое отчуждение кадров.
    \item Согласно методологии академика Ж. Т. Тощенко, современная российская НТИ переживает глубокую внутреннюю стратификационную поляризацию \cite{Toschenko}. Наряду с успешной, интегрированной во власть технобюрократией госкорпораций и конформным массовым слоем ИТР, фиксируется появление «интеллектуального прекариата» --- ученых и инженеров гражданского сектора, находящихся в состоянии социально-экономической нестабильности и ценностной аномии, что порождает скрытый протестный потенциал.
    \item Протестная активность НТИ носит специфический технократический характер, выражаясь не в уличных манифестациях, а в экспертной критике некомпетентности аппарата, уходе в приватную сферу и деструктивном для государства феномене «утечки мозгов». Преодоление этих кризисных явлений и обеспечение технологического суверенитета РФ требуют от власти перехода к стратегии гибкого партнерства, основанной на предоставлении НТИ реальной академической автономии, преодолении прекаризации труда инженеров и открытии прозрачных каналов вертикальной мобильности в структуре управления государством.
\end{enumerate}